\chapter{Projekt architektury systemu}
\label{chap:architecture}
\section{Podział odpowiedzialności}
Rysunek~\ref{fig:architecture} przedstawia architekturę docelową. W pierwszym etapie broker i aplikacja bramy mają działać na laptopie, a rolę urządzeń pełni symulator.
\begin{figure}[htbp]
\centering
\begin{tikzpicture}[
  box/.style={draw=draftblue,rounded corners=2pt,align=center,
    text width=4.9cm,minimum height=1.05cm,font=\small},
  arrow/.style={-{Latex[length=2mm]},thick,draw=draftblue},
  node distance=0.6cm and 0.7cm]
\node[box] (sim) {Symulator na laptopie\\konfigurowalne węzły};
\node[box,right=of sim] (esp) {Fizyczne węzły ESP32\\dołączane później};
\node[box,below=of sim] (mqtt) {Broker MQTT};
\node[box,below=of mqtt] (edge) {Kolektor i walidacja\\SQLite oraz FastAPI};
\node[box,right=of edge] (model) {Trening na laptopie\\eksport modelu};
\node[box,below=of edge] (infer) {Późniejsza inferencja na bramie\\anomalie i alerty};
\draw[arrow] (sim) -- (mqtt);
\draw[arrow] (esp.south) |- (mqtt.east);
\draw[arrow] (mqtt) -- (edge);
\draw[arrow] (edge) -- (infer);
\draw[arrow] (model.south) |- (infer.east);
\end{tikzpicture}

\caption{Planowany przepływ danych w EdgeGuard IoT. Schemat nie potwierdza ukończenia komponentów.}
\label{fig:architecture}
\end{figure}

\section{Identyfikacja urządzeń i komponentów}
\writingnote{Opisz device\_id, identyfikatory komponentów, role, strefy i możliwości urządzeń. Pokaż konfiguracje z różnym podziałem czujników, bez utrwalania GPIO w aplikacji bramy.}

\section{Kontrakty MQTT i walidacja danych}
\writingnote{Udokumentuj tematy, wersję schematu, czas urządzenia i odbioru, identyfikator uruchomienia, numer sekwencyjny, brakujące pomiary oraz obsługę duplikatów.}

\section{Przechowywanie danych i API}
\writingnote{Dodaj diagram bazy oraz opis encji i indeksów po ich implementacji. Powiąż przykłady odpowiedzi API z rzeczywistymi testami.}

\section{Integracja detektora i obsługa modeli}
\writingnote{Przedstaw interfejs AnomalyDetector, zgodność cech, wersjonowanie i wycofanie modelu. Oddziel trening na laptopie od inferencji na bramie.}

\section{Niezawodność i bezpieczeństwo}
\writingnote{Uzasadnij kolejki, limity, obsługę ponownych połączeń, statusy oraz późniejsze dane dostępowe i uprawnienia MQTT.}
